%% Consumer-corrected derivative of the historical A-162 notation tail.
\section{Model-theory notation lookup}
\label{sec:p03-model-symbols}

Table~\ref{tab:p03-model-symbols} is a separately authored lookup of the
registry's existing model-theory assignments.  In particular,
\(\cong\), never the Kleene-equality symbol reserved for Part~4, denotes
isomorphism here.
\index{elementary equivalence}

\begingroup\small
\begin{longtable}{@{}p{0.29\textwidth}p{0.61\textwidth}@{}}
\caption{Model-theory symbol subtable.}
\label{tab:p03-model-symbols}\\
\toprule
Notation & Registered meaning \\
\midrule
\endfirsthead
\toprule
Notation & Registered meaning \\
\midrule
\endhead
\(\mathcal M\preccurlyeq\mathcal N\) &
  \(\mathcal M\) is an elementary substructure of \(\mathcal N\). \\
\addlinespace
\(\mathcal M\equiv\mathcal N\) &
  \(\mathcal M\) and \(\mathcal N\) satisfy exactly the same
  \(\Lang\)-sentences (elementary equivalence). \\
\addlinespace
\(\operatorname{Diag}(\mathcal M)\) &
  Atomic and negated atomic \(\Lang(M)\)-sentences true in \(\mathcal M_M\). \\
\addlinespace
\(\operatorname{ElDiag}(\mathcal M)\) &
  All \(\Lang(M)\)-sentences true in \(\mathcal M_M\). \\
\addlinespace
\(S_n(T)\) &
  Space of complete \(n\)-types over the empty set consistent with \(T\), with
  its Stone topology; see Definition~\ref{def:type-space}. \\
\addlinespace
\(\displaystyle\prod_{i\in I}\mathcal M_i/\mathcal U\) &
  Ultraproduct of the family \(\langle\mathcal M_i:i\in I\rangle\) modulo
  the ultrafilter \(\mathcal U\); see
  Sections~\ref{sec:types-ultraproducts-04}
  and~\ref{sec:types-ultraproducts-05}. \\
\addlinespace
\(\mathcal M\cong\mathcal N\) &
  \(\mathcal M\) and \(\mathcal N\) are isomorphic. \\
\bottomrule
\end{longtable}
\endgroup
